\subsection{Deterministic wealth and the Dynkin integrand}
Throughout this section the primitives, controls, and restart are those stated immediately above. With $p=0$, $c=3/4$, and $x(0)=5/4$,
\begin{equation}
 x(t)=\frac{75}{2}-\frac{145}{4}e^{t/50}.
\end{equation}
For $0\le t\le1$, $e^{t/50}\le(1-1/50)^{-1}=50/49$, hence $x(t)\ge25/49>1/2$. It is decreasing and starts below two, so no wealth boundary is reached. Let $y$ denote preference. Direct differentiation of the unchanged settlement gives
\begin{align}
 h_\theta(t,y)={}&\frac{(3/4)^{1-y}}{1-y}-\theta^2+8
 -.04\theta(y-2)-.00005+.002-\frac{.075}{x(t)}\nonumber\\
 &\quad+.0008(y-2)^2-.004\log x(t)+.32(1-t),\label{eq:hsupp}\\
 \partial_\theta h_\theta={}&-2\theta-.04(y-2).\label{eq:hparametersupp}
\end{align}
For $y\in[6/5,14/5]$, put $z=y-1\in[1/5,9/5]$ and $L=\log(4/3)<1/3$. The utility is $-e^{Lz}/z$, which is increasing on this interval because $Lz<1$. Its lower bound is $-5e^{1/15}\ge-75/14$. Also $|y-2|\le4/5$, $x(t)\ge1/2$, and $\log x(t)\le\log(5/4)<1/4$. Dropping nonnegative terms yields
\begin{equation}
 h_\theta\ge8-\frac1{20000}-\frac3{20}-\frac1{1000}
 -\frac4{625}-\frac1{25}-\frac{75}{14}
 =\frac{342357}{140000}>2.\label{eq:hlower}
\end{equation}
The utility derivative is $e^{Lz}(1-Lz)/z^2\ge10/81$. Thus
\begin{equation}
 \partial_yh_\theta\ge\frac{10}{81}-\frac4{3125}-\frac1{125}
 =\frac{28901}{253125}>0.\label{eq:hylower}
\end{equation}
Extend $h_\theta(t,y)$ constantly from $b=14/5$ to larger $y$. The extension remains bounded, nondecreasing in $y$, and at least two. Its parameter derivative retains the bounds in \eqref{eq:hparametersupp}. Conservative absolute bounds $|h|\le20$ and $|\partial_\theta h|\le2$ hold for this extension.

\subsection{A coupling lower bound for the derivative}
Let $Y_t=u_0+\theta t+\sigma W_t$, killed at $a=6/5$, where $u_0=61/50$ and $\sigma=1/20$. If drift increases, the coupled process is higher at every date and lower-boundary survival events are nested. For a fixed nondecreasing integrand $h\ge2$, paths surviving under both drifts have a nonnegative payoff change, while newly surviving paths contribute at least two. Taking a positive difference quotient gives
\begin{equation}
 \partial_\theta\E_\theta[h(t,Y_t)\ind\{t<\tau_a\}]
 \ge2\partial_\theta\Pr_\theta(t<\tau_a).
\end{equation}
The direct dependence of $h_\theta$ on $\theta$ must be added, not included in this frozen-integrand comparison. It contributes at least $-4/125$ for $\theta\in[-1/5,0]$ and at least $-54/125$ for $\theta\in[0,1/5]$, after integrating over a discounted survival time of at most one. Gaussian likelihood scores justify differentiability and dominated passage to the occupation integral for this bounded extension.

The reflection formula with drift, differentiated exactly, yields
\begin{equation}\label{eq:survderiv}
 S'_\theta(t)=16e^{-16\theta}\Phi\left(\frac{\theta t-1/50}{(1/20)\sqrt t}\right).
\end{equation}
The function $\log\Phi$ is concave. Indeed, with $r=\varphi/\Phi$, its second derivative is $-r(z)\{z+r(z)\}<0$; when $z<0$, the elementary Gaussian tail inequality gives $r(z)>-z$. Equation~\eqref{eq:survderiv} is therefore log-concave in $\theta$. Its minimum on each of the two half intervals is at least the smaller of its endpoint values.

For $t\in[1/25,1/5]$, the normal arguments at $\theta=0$ and $\theta=-1/5$ are bounded below by $-2$ and $-14/5$, respectively. The latter follows by maximizing $4\sqrt t+(2/5)/\sqrt t$ on that compact interval; its maximum is $14/5$. The rational inequalities
\begin{gather}\label{eq:normalbounds}
 24<e^{16/5}<25,\qquad \Phi(-2)>1/50,\qquad \Phi(-14/5)>1/500,\\
 \Phi(-4/5)>1/5,\qquad \Phi(6/5)>22/25.\nonumber
\end{gather}
imply $S'_\theta(t)\ge8/25$ on the entire negative drift half interval.

For $t\in[1/4,1]$, the endpoint arguments at $\theta=0$ and $\theta=1/5$ are at least $-4/5$ and $6/5$. Hence $S'_\theta(t)\ge352/625$ on the positive drift half interval. Since $e^{-.04t}\ge24/25$, the two lower bounds for the derivative of the lower-killed occupation expression are
\begin{align}
 2\frac4{25}\frac{24}{25}\frac8{25}-\frac4{125}&=.066304,\\
 2\frac34\frac{24}{25}\frac{352}{625}-\frac{54}{125}&=.379008.
\end{align}
No derivative sample grid appears in this argument. The intervals cover every admissible constant drift, including the shared point zero.

\subsection{Distant-boundary error and rational checks}
The upper boundary is distant but not discarded. Under all admissible drifts, a path hitting it must have a Brownian excursion of at least $(b-u_0-1/5)/\sigma$. Reflection and a Gaussian tail bound give
\begin{equation}
 p_b\le2\exp\left[-\frac{(b-u_0-1/5)^2}{2\sigma^2}\right]
 =2e^{-380.88}<2^{-379}.
\end{equation}
For a path-dependent difference supported on that event, the drift likelihood score has second moment at most $1/\sigma^2$. Cauchy--Schwarz bounds the derivative discrepancy between the actual doubly killed occupation integral and the extended lower-killed one by
\begin{equation}
 \frac{20}{\sigma}\sqrt{p_b}+2p_b
 <400\,2^{-189}+2\,2^{-379}<1/1000.
\end{equation}
Subtracting $1/1000$ from the preceding half-line bounds gives precisely $8163/125000$ and $47251/125000$. A strictly increasing continuous function on the compact drift interval has the unique maximizer $1/5$.

The file \path{revisions/2026-09-24-r34/replication/boundary.py} proves the elementary constants in \eqref{eq:normalbounds} with exact fractions. Exponential bounds use Taylor sums and an explicit geometric tail. A Machin arctangent identity with alternating-series bounds proves a rational bound for $\pi$, and $\sqrt{2\pi}<251/100$ follows. Gaussian lower bounds use monotone right sums of mesh $1/100$ on finite intervals, with every exponential term bounded by a rational Taylor calculation. Truncating the infinite tail only lowers a tail integral. For $\Phi(6/5)$ the integral is from zero to $6/5$, added to $1/2$. Terms are rounded downward to a common dyadic denominator before summation, with the direction of the bound preserved. The code records the resulting exact rational lower bounds and requires them to exceed all displayed thresholds.

The payoff enclosure at $\theta=1/5$ is inherited from the unchanged validated R30 killed-kernel oracle. Its source path, Git blob hash, and the fact that it is not recomputed are recorded in \path{revisions/2026-09-24-r34/results/boundary.json}. The current theorem changes its interpretation from an endpoint evaluation to the optimum of the entire constant class. It does not assert feedback optimality or satisfy the original all-domain current-state target.

